% =====================================================================
% Stage 1 Concept and Feasibility Report -- Part B
% Tool return confirmation at aircraft maintenance sign-off
% Project I (247-519-VA)
%
% Draft prepared with an AI coach from the student's own sessions, logbook
% numbers and decisions. Every \todo{} (red) is a value only the student has
% (logbook page, date, weight, count) or a result not measured yet.
% Nothing marked \todo{} has been guessed. Every \inferred{} (blue) was filled
% on Nov 18 from session data or chat notes and must be confirmed.
%
% Standalone: compiles as is. To merge into main.tex, copy everything
% between \begin{document} and \end{document}.
% =====================================================================
\documentclass[conference]{IEEEtran}
\usepackage[T1]{fontenc}
\usepackage{booktabs}
\usepackage{tabularx}
\usepackage{makecell}
\usepackage{xcolor}
\usepackage{amsmath}

\providecommand{\todo}[1]{\textcolor{red}{[#1]}}
\newcolumntype{Y}{>{\raggedright\arraybackslash}X}
\renewcommand{\arraystretch}{1.15}

\begin{document}

\title{Stage 1 Concept and Feasibility Report -- Part B\\
\large Tool Return Confirmation at Aircraft Maintenance Sign-off}
\author{\IEEEauthorblockN{Gib Pia Macarine}
\IEEEauthorblockA{Project I (247-519-VA), Vanier College\\ November 18, 2026}}
\maketitle

\begin{abstract}
This report records what changed since the frozen proposal, compares the
candidate concepts, and sets out the evidence behind each claim. Bench
sessions on a shadow board showed that a before/after (checkout versus
sign-off) comparison detects removed tools and runs well under the 3\,s
limit, but misses some look-alike swaps and gives false alarms on a complete
kit. The chosen concept is therefore a hybrid: the before/after comparison
finds which slots changed, and a learned slot classifier names what is in
each changed slot. The recommendation is to proceed with conditions; each
open claim has a planned test and a date before Gate~2.
\end{abstract}

% ---------------------------------------------------------------------
\section{Change Log}
% ---------------------------------------------------------------------
Identifiers from the frozen proposal (N1--N7, REQ-TR-01 to 08, R1--R3,
A1--A2, D1, C1--C4) are kept and never reused. Table~\ref{tab:changes} lists
every change since the proposal, with its reason and source.

\begin{table*}[!t]
\centering
\caption{Changes since the frozen proposal}
\label{tab:changes}
\begin{tabularx}{\textwidth}{@{}l l Y Y l@{}}
\toprule
\# & Date & Change & Reason & Source \\
\midrule
1 & Oct 16 & Custom hardware runs inference; computer only trains and logs; watch is the indicator only. & Instructor feedback: the watch should not run the model. & \inferred{L-08} \\
2 & Nov 8 & Kit cut from 16 to 12 tools; all five look-alikes kept. & Not all 16 fit in frame; edge tools touched the border. & L12 \\
3 & Nov 9 & Blank paper replaced by a shadow board (pink tool-shaped shadows, traced outlines). & Dry run L12: free placement on blank paper failed for both methods. & L12, \inferred{L13} \\
4 & \todo{date} & Claims numbered CL1, CL2, \ldots & The proposal already uses C1--C4 for platform constraints. & \todo{L-xx} \\
5 & \inferred{Nov 12} & Lighting is not built yet; a retake check is used as the lighting mitigation for now. & Lighting listed in the proposal; not available for the PoC. & \inferred{L13} \\
6 & Nov 13 & Prototype processor left open; the Pi Zero~2~W is used only as a measurement stand-in. & Memory and compute needs must be measured before choosing. & \inferred{L14} \\
7 & Nov 17 & Recognition changed to a hybrid: before/after comparison plus a learned slot classifier. & Method A alone cannot name the tool in a slot and is not a learned model, so it does not meet the AI core-domain rule. & \inferred{L15} \\
8 & Nov 17 & Retake check added; colour floor lowered from 6.0 to 0.1; camera exposure locked (now $-4$). & Shadow trials caused false alarms (Nov 12 session); the old colour floor sat above every observed colour score. & L15 \\
\bottomrule
\end{tabularx}
\end{table*}

% ---------------------------------------------------------------------
\section{Concept Alternatives}
% ---------------------------------------------------------------------
Five concepts were considered. Each fails for a different reason, which is
why they are worth comparing.

\begin{itemize}
\item \textbf{A -- Before/after comparison.} At checkout and at sign-off the
camera photographs the board; each slot is compared with the same job's
checkout photo. It fails on look-alike swaps and cannot see a tool that was
already missing at checkout.
\item \textbf{B -- Known shadow board.} Each slot is compared with reference
photos stored once at setup. It fails by drifting: once a tool is put back
slightly differently, that slot no longer matches its reference (the pliers
in session~2 stayed above threshold for every trial after being reseated).
\item \textbf{C -- Manual check.} The technician or storeman counts the kit by
eye. This is the status quo, and it fails through human attention.
\item \textbf{D -- Hybrid (A + learned slot classifier).} A flags which slots
changed; a classifier trained on the computer names what is in each changed
slot: its own tool, empty, or another tool. Running the classifier on the
checkout photo also covers A's blind spot.
\item \textbf{E -- Fiducial tags on each tool.} A printed tag (e.g.\
AprilTag) on every tool identifies it when the tag is seen. Not scored:
at the bench geometry a 1\,cm tag spans about 18\,px, under 2\,px per tag
cell (calculated), below what tag detectors read reliably. Tags also wear
off or get covered, and depending on every tool being tagged is the weakness
of the existing tagging systems the proposal set out to avoid. Tag decoding
is not a learned model, so E would not make AI a core domain.
\end{itemize}

Concepts A, B and C were scored against six weighted criteria (Table~\ref{tab:scores}).
D was added after scoring; it inherits A's measured strengths, and its
claimed improvement (swaps) is tested directly as CL2.

\begin{table}[!t]
\centering
\caption{Weighted concept scores (1 = poor, 5 = good)}
\label{tab:scores}
\begin{tabular}{@{}lcccc@{}}
\toprule
Criterion & Weight & A & B & C \\
\midrule
Buildability          & \todo{w} & \todo{} & \todo{} & \todo{} \\
Cost                  & \todo{w} & \todo{} & \todo{} & \todo{} \\
Timing                & \todo{w} & \todo{} & \todo{} & \todo{} \\
Misses                & \todo{w} & \todo{} & \todo{} & \todo{} \\
False alarms          & \todo{w} & \todo{} & \todo{} & \todo{} \\
Remaining uncertainty & \todo{w} & \todo{} & \todo{} & \todo{} \\
\midrule
\textbf{Weighted total} & 1.00 & \textbf{3.60} & \textbf{2.45} & \textbf{3.20} \\
\bottomrule
\end{tabular}
\end{table}

B is rejected on measured drift (CL6, refuted). C is rejected because it
gives no record and depends entirely on attention, the problem the project
exists to address. D is chosen because it keeps A's measured removal and
timing results, targets A's measured weakness, and makes AI a core domain:
a learned model with a measured error rate.

\begin{figure}[!t]
\centering
\fbox{\parbox{0.9\columnwidth}{\centering\vspace{1.2cm}
\todo{Block diagram of concept D (drawn by the student)}\vspace{1.2cm}}}
\caption{Block diagram of the chosen concept (D, hybrid).}
\label{fig:block}
\end{figure}

% ---------------------------------------------------------------------
\section{Feasibility Case}
% ---------------------------------------------------------------------
Each claim tests a requirement or risk, and each is labelled
\emph{measured} (run, with conditions stated), \emph{calculated} (worked out
from a stated model) or \emph{assumed} (believed, with what would settle it).

\subsection{Evidence so far}
\textbf{Dry run (L12, Nov 8).} On blank paper with 16 tools, A missed 5/10
removals and B 7/10; both missed all 10 swaps. It is kept as evidence that
free placement does not work; no thresholds were taken from it.

\textbf{Session 2 (Nov 12, shadow board, 12 tools).} Method A caught 12/12
removals and 8/10 look-alike swaps, with 4 false alarms on complete-kit
captures (out of 10). Every check took at most 1.13\,s on the PC.

\textbf{Session 3 (Nov 17, L15, retake check active).} Thresholds were
recalibrated and frozen before the trials (12 complete-kit captures,
$\max(\text{floor}, \text{mean} + 4\sigma)$). Table~\ref{tab:l15} gives the
results; all 32 trials were valid.

\begin{table}[!t]
\centering
\caption{Session 3 results (L15), 95\,\% one-sided upper bounds}
\label{tab:l15}
\begin{tabular}{@{}lcc@{}}
\toprule
 & Method A & Method B \\
\midrule
Removals caught      & 11/12 & 11/12 \\
Swaps detected       & 10/10 & 7/10 \\
Complete-kit false alarms & 0/10 ($<$25.9\,\%) & 0/10 \\
Slowest check (s)    & 1.29 & 1.21 \\
\bottomrule
\end{tabular}
\end{table}

\textbf{The one miss came from calibration, not detection.} During
calibration, one replacement of the wire cutter scored 0.27, as high as a
real removal, and alone raised that slot's threshold to 0.33. Without that
capture the threshold would have been about 0.10, and the trial-11 removal
(0.26) would have been flagged. The miss is reported as measured, since the
thresholds were frozen first. From now on, calibration scores are reviewed
before freezing, and a capture where a lifted tool scores like a removal is
redone.

\textbf{Swaps are detected but not named.} Method A labels a slot only as
``changed''. It caught the cutter swaps through colour, not shape, so
naming the tool in a slot still depends on the classifier (CL2).

\textbf{Two unplanned findings.} In trial 14, a swap from trial 13 had not
been undone. Method A treated the wrong placement as normal, because its
checkout photo already contained it, while Method B flagged it. This is
risk R7 observed directly. Later, after the nippers were reseated in trial
16, Method B kept flagging them although nothing was wrong: the same drift
seen with the pliers in session 2. Trial 16 was an operator error (trial 15
not restored) and is not counted as a system false alarm.

\textbf{Reading the numbers.} With few trials, a clean result proves little:
0 misses in 12 removals only bounds the miss rate below 22\,\% (one-sided
95\,\%), and 1 in 12 gives an upper bound of 34\,\%. The planned 90 removals
with 0 misses would bring the bound to 3.3\,\%, which is why the full run
matters.

\subsection{Claims}
Table~\ref{tab:claims} lists the claims. CL1--CL5 all come from the same bench
captures, so one session tests them together: every check records which
slots were changed, what the system reported, and how long it took.

\begin{table*}[!t]
\centering
\caption{Claims, what they test, and their status}
\label{tab:claims}
\begin{tabularx}{\textwidth}{@{}l Y l Y l@{}}
\toprule
ID & Claim & Tests & Evidence so far & Status \\
\midrule
CL1 & Every removed tool is flagged as missing. & REQ-TR-02, R1 & 12/12 (S2); 11/12 (S3, miss traced to calibration). Plan: 90 removals, pass at 0 misses. & Measured (partial) \\
CL2 & A look-alike swap is flagged and the tool in the slot is named. & REQ-TR-02, REQ-TR-01, R1 & A alone: 8/10 (S2), 10/10 (S3) detected, none named. Hybrid: not run yet. & Detection measured; naming open \\
CL3 & A complete kit is not flagged. & REQ-TR-06 & 4/10 (S2); 0/10 (S3, with retake). Plan: 300 checks, pass at 0. & Measured (partial) \\
CL4 & A check, capture to result, takes 3\,s or less. & REQ-TR-07, R3 & At most 1.13\,s (S2) and 1.29\,s (S3) on the PC. & Measured (stand-in) \\
CL5 & A shadow or light change triggers a retake instead of a wrong verdict. & R2 & S3: retake on 3/3 lighting trials (26, 28, 32); trial 26's first capture would have flagged all 12 slots. Unneeded retakes: \todo{count}. & Measured (partial) \\
CL6 & A stored reference (Method B) stays valid over repeated returns. & Concept B & Pliers (S2) and nippers (S3) drifted above threshold after one reseat. & Refuted \\
CL7 & The check fits the memory and compute of a low-cost processor. & REQ-TR-05, R4 & None yet. Plan: replay stored images on the Pi Zero~2~W. & Open \\
CL8 & The watch shows the result; a lost link shows NOT CONFIRMED within 10\,s. & REQ-TR-04, REQ-TR-01 & Sketch ready, not run. Plan: 60/60 with the link cut. & Open \\
CL9 & The custom hardware costs \$75 or less, excluding watch and computer. & REQ-TR-03 & Parts list started; prices still to be filled. & Open \\
\bottomrule
\end{tabularx}
\end{table*}

The record requirement (REQ-TR-08) is covered by the same bench session:
every check is written as one row of a CSV file.

\subsection{Data for the learned classifier}
Only one physical copy of each tool exists, and that matches the real use: a
shadow board holds one specific kit, so the model must recognize \emph{this}
kit's tools in \emph{this} board's slots, not tools in general. Variety comes
from conditions, not copies: placement, lighting and swaps. Every stored
capture gives 12 labelled slot crops, labelled from the trial plan rather
than by eye.

\begin{itemize}
\item The data are split by session, never by image, so the test session
holds conditions the model has not seen.
\item Swap crops for the two tapes and for the cutters are treated with
caution: where the pink shadow shows, the model may learn the shadow rather
than the tool. They are tagged and reported separately.
\item Calibration photos with visible misalignment (e.g.\ scissors showing
pink) are reviewed before use. The shadow fixes each tool's orientation, so
this is expected to be rare.
\end{itemize}

\subsection{Test plan to close the open claims}
The work is done in order, each step on its own git branch, so that a
stable version is committed before the next step builds on it.
\begin{enumerate}
\item \textbf{Retake check (CL5).} Done in session 3 and committed as the
stable base. Still to do: log retakes in \texttt{trials.csv} (the columns were
missing in session 3, so retakes were recovered from saved images and notes),
and report lighting retakes apart from ``many slots flagged'' retakes.
\item \textbf{Hybrid model (CL2).} Train the slot classifier on stored crops,
evaluate it with Method A on a held-out session, then run the full tests:
90 removals (CL1), 300 complete-kit checks (CL3), with timing (CL4).
\item \textbf{Watch link (CL8, CL4).} Simulate first, then connect the PC;
the 60 checks with the link cut test the timeout. Pressing \emph{check} on the
watch can replay a stored capture, giving end-to-end latency without the
live camera. Replay leaves out capture time, which is reported separately.
\end{enumerate}

% ---------------------------------------------------------------------
\section{Baselined Requirements}
% ---------------------------------------------------------------------

\begin{table*}[!t]
\centering
\caption{Requirements after Stage 1}
\label{tab:reqs}
\begin{tabularx}{\textwidth}{@{}l Y Y l l@{}}
\toprule
ID & Requirement (short form) & Pass criterion & Claim & Status \\
\midrule
REQ-TR-01 & The watch names each missing tool. & With 2 tools removed, the watch names both & CL2, CL8 & Partial \\
REQ-TR-02 & Misses fewer than 1 in 30 missing tools. & 0 misses in 90 removals & CL1 & Partial \\
REQ-TR-03 & Custom hardware cost limit. & \$75 or less (\todo{CAD/USD}) & CL9 & Open \\
REQ-TR-04 & No result within 10\,s reads as not confirmed. & 60/60 NOT CONFIRMED with link cut & CL8 & Open \\
REQ-TR-05 & A model trained on the computer runs on the custom hardware. & A model trained on the computer classifies a kit image on the custom hardware & CL7 & Open \\
REQ-TR-06 & False alarm on fewer than 1 in 100 complete kits. & 0 false alarms in 300 checks & CL3 & Partial \\
REQ-TR-07 & Check completes in 3\,s or less. & 90 timed checks, none over 3\,s & CL4 & Stand-in \\
REQ-TR-08 & The computer stores a dated record of every check, listing each tool present. & After 10 checks, the record lists the time and tools of each & bench CSV & Partial \\
\bottomrule
\end{tabularx}
\end{table*}

REQ-TR-05 keeps its proposal wording: the hybrid's classifier is trained on
the computer and exported to run on the processor, so the change of
recognition method does not change this requirement.

% ---------------------------------------------------------------------
\section{Risk Register}
% ---------------------------------------------------------------------
\noindent Likelihood (L) and impact (I) are ranked from 1 (low) to 5 (high).

\begin{table*}[!t]
\centering
\caption{Ranked risks}
\label{tab:risks}
\begin{tabularx}{\textwidth}{@{}l Y c c Y@{}}
\toprule
ID & Risk & L & I & Mitigation and evidence \\
\midrule
R1 & Look-alike tools (stripper, cutter, nippers; two tapes) are swapped and not caught. & \inferred{4} & \inferred{5} & Hybrid classifier names slot content (CL2). \\
R2 & Shadows or light changes cause false alarms or misses. & \inferred{4} & \inferred{4} & Retake check now (CL5, 3/3 in S3). It handles temporary changes only; a lasting change uses up the retakes and then flags. Built lighting plus retake at Gate 2. \\
R3 & End-to-end time through the watch exceeds what a technician will wait. & \inferred{2} & \inferred{3} & Watch replay test (CL4, CL8). \\
R4 & The check does not fit a low-cost processor's memory or compute. & \inferred{3} & \inferred{4} & Pi replay profile at three resolutions (CL7). \\
R5 & Training on one copy per tool does not generalize to new conditions. & \inferred{3} & \inferred{4} & Session-based split; vary lighting and placement between sessions. \\
R6 & The model learns the pink shadow instead of the tool. & \inferred{3} & \inferred{3} & Tag and report those crops separately. \\
R7 & A tool already missing or misplaced at checkout is treated as normal by A. & \inferred{3} & \inferred{5} & Observed in S3, trial 14. Run the classifier on the checkout photo (concept D). \\
R8 & One bad calibration capture inflates a slot's threshold and causes a miss. & \inferred{3} & \inferred{4} & Observed in S3 (wire cutter). Review calibration scores before freezing; redo outlier captures. \\
R9 & Nuisance alarms make operators stop trusting the check. & \inferred{3} & \inferred{3} & Observed with B's drift in S3. B rejected; false alarms tracked (CL3). \\
\bottomrule
\end{tabularx}
\end{table*}

Method B's drift is no longer a project risk, since B was rejected.

% ---------------------------------------------------------------------
\section{Gate 1 Readiness}
% ---------------------------------------------------------------------
\textbf{Recommendation: proceed with conditions.}

\textbf{Why.} Over two sessions Method A caught 23/24 removals, its one miss
traced to a calibration capture, and detected 18/20 swaps, all 10 in the
latest session. Every check took under 1.3\,s against the 3\,s limit. The
retake check turned all three lighting disturbances into retakes instead of
false alarms. Method B is rejected on drift measured twice. What remains is
naming the tool in a changed slot. That has a planned fix (the hybrid) and a
test that can come out as ``no''.

\textbf{Conditions before Gate 2}, in order:
\begin{enumerate}
\item CL5: retake columns logged; unneeded-retake rate counted -- \todo{date}
\item CL2 hybrid evaluated on a held-out session -- \todo{date}
\item CL1 / CL3 / CL4: 90 removals and 300 complete-kit checks -- \todo{date}
\item CL8 watch link and 60-check timeout -- \todo{date}
\item CL7 Pi replay profile; CL9 priced parts list -- left open, \todo{date}
\end{enumerate}

\textbf{Beyond the PoC}, attempted if time allows, in order of priority:
a Pi Zero~2~W run replaying stored captures (and a build attempt), wider
lighting conditions, and a dashboard for the storemen or a supervisor.

\textbf{What Gate 2 will show.} A prototype joining camera, processor and
watch, running the hybrid check with both lighting mitigations combined:
the retake check and built lighting.

\end{document}
